%---------------------------------------------------------------
%  SECTION I — INTRODUCTION
%---------------------------------------------------------------
\section{Introduction}

The global proliferation of cloud storage services, remote work
infrastructure, and inter-organisational data exchange has made
secure file transfer a critical requirement. Current industry
deployments of Secure Shell File Transfer Protocol (SFTP)
\cite{ylonen2006ssh} rely on RSA-2048 or ECDH for key agreement
and RSA or ECDSA for authentication. Shor's algorithm
\cite{shor1994algorithms}, when executed on a cryptographically
relevant quantum computer, solves the integer factorisation and
discrete logarithm problems in polynomial time, directly breaking
all of these classical schemes.

The threat is not merely theoretical. The ``Harvest Now, Decrypt
Later'' (HNDL) adversarial model \cite{mosca2018cybersecurity}
posits that an adversary can capture encrypted traffic today and
decrypt it retroactively once quantum hardware matures. Data with
long-term sensitivity (government records, medical files, financial
contracts) is therefore at risk \emph{now}. Grover's algorithm
\cite{grover1996fast} additionally provides a quadratic speedup for
symmetric key search, halving the effective security of AES-128 to
64~bits against a quantum attacker.

The National Institute of Standards and Technology (NIST) concluded
its Post-Quantum Cryptography (PQC) standardisation process in 2024,
publishing FIPS~203 (ML-KEM / Kyber) \cite{nist2024pqc},
FIPS~204 (ML-DSA / Dilithium), and FIPS~205 (SLH-DSA). These
lattice-based algorithms are believed to resist both classical and
quantum adversaries under the Module Learning With Errors (MLWE)
assumption.

Despite the urgency, PQC adoption in production file transfer systems
remains scarce. Most existing works focus on isolated cryptographic
primitives or theoretical protocol analyses \cite{dam2023survey},
leaving a gap for end-to-end, deployable implementations.
This paper addresses that gap with \textbf{Q-SFTP}, which contributes:

\begin{itemize}
  \item A four-phase post-quantum mutual authentication handshake
        (Q-TLS) using Kyber512 and Dilithium2.
  \item A BLAKE2b-256 \cite{aumasson2013blake2} dual-side integrity
        ecosystem with persistent hash registry and background monitoring.
  \item A chunked resumable transfer protocol with cross-session
        state persistence via JSON log and browser \texttt{localStorage}.
  \item Defense-in-depth: RBAC, metadata stripping, malicious-file
        detection, and tamper-evident audit logging.
  \item A browser-based web GUI and Docker containerisation for
        production deployment.
\end{itemize}

The remainder of this paper is organised as follows.
Section~\ref{sec:background} reviews related work and cryptographic
background. Section~\ref{sec:architecture} details the system design
and protocol. Section~\ref{sec:results} presents experimental results
and security analysis. Section~\ref{sec:conclusion} concludes with
future directions.
